\documentclass[english, a4paper, 12pt, twoside, openany]{book}

% -------------- Setup, no cambiar ---------------
% Codificación, idioma y tipografía
\usepackage[T1]{fontenc}                         % Para las fuentes (T1 encoding)
\usepackage[utf8]{inputenc}                      % Introduce la codificación utf8 que permite el uso de tildes y otros.
\usepackage{textcomp}                            % Agrega muchos símbolos extra para utilizar en el texto
\usepackage[english]{babel}                      % Idioma
\usepackage{microtype}                           % Permite la modificación de los caracteres.

% Configuración de página y espaciado
\usepackage[left=4cm,top=4cm,bottom=2.5cm,right=2.5cm]{geometry}    % Márgenes del documento
\usepackage{setspace}                                               % Permite elegir el interlineado
\setstretch{1.3}                                                    % Interlineado de uno y medio. 1.6 es interlineado doble.
\usepackage{parskip}                                                % Elimina la sangría
\usepackage{titlesec}                                               % Suplemento a fancyhdr
\setcounter{secnumdepth}{4}                                         % El contador de figuras, secciones, tablas, etc. llega hasta profundidad de 4. i.e. 1.1.1.1
\hyphenpenalty=100000                                               % No se dividen las palabras al terminar una línea
\sloppy                                                             % Ayuda a que las cosas no se salgan de los márgenes.
\raggedbottom                                                       % Hace que el tamaño de las páginas llegue hasta donde llegue el texto, no agrega espacio vertical

% Paquetes para gráficos, figuras y tablas
\usepackage[table,xcdraw,dvipsnames]{xcolor}    % Agranda la cantidad de colores.
\usepackage{graphicx}                           % Para poner imágenes
\graphicspath{{Images/}{../Images/}}            % Dónde estarán las imágenes
\usepackage{adjustbox}                          % Permite más versatilidad para cosas como includegraphics.
\usepackage{float}                              % Permite poner figuras en cualquier lugar.
\usepackage{wrapfig}                            % Permite poner figuras con texto alrededor.
\usepackage[rightcaption]{sidecap}              % Permite agregar texto lateral
\usepackage[labelfont=bf]{caption}              % Texto en negrita para descripciones (\caption)
\usepackage{subcaption}                         % Permite subtextos en las figuras. Es incompatible con el paquete subfig y subfigure
\usepackage{array}                              % Para ordenar texto y ecuaciones.
\usepackage{multirow}                           % Extensión para tablas
\usepackage{makecell}                           % Ayuda en la creación de tablas
\usepackage{hhline}                             % Agranda las opciones de las líneas
\usepackage[para]{threeparttable}               % Tablas vistosas, mirar antes de utilizar.

% Matemáticas y numeración
\usepackage{amsmath, amssymb, amsthm}           % Paquetes matemáticos
\numberwithin{equation}{section}                % Numeración de ecuaciones
\usepackage{chngcntr}                           % Necesario para enumerar las tablas
\counterwithin{table}{section}                  % Numeración de tablas
\counterwithin{figure}{section}                 % Numeración de figuras

% Organización del documento y referencias internas
\usepackage{subfiles}                           % Para agregar los capítulos que se escriben aparte.
\usepackage[toc,page]{appendix}                 % Permite el uso de apéndices
\usepackage{xparse}                             % Ayuda en la creación de nuevos comandos
\usepackage{nameref}                            % Referencias por nombre
\usepackage[bottom,hang,flushmargin]{footmisc}  % Notas al pie de página están fijas al fondo y sin sangría
\interfootnotelinepenalty=10000                 % Previene que las notas al pie de página se pasen a la siguiente página cuando son muy largas
\usepackage{fancyhdr}                           % Permite formatear las cabeceras, pies, enumeración, etc.

% Hipervínculos y URLs
\usepackage{url}                                % Permite el uso de enlaces URL.

% Contenido auxiliar
\usepackage{lipsum}                             % Para crear texto de relleno.

% Referencias bibliográficas
\usepackage{natbib}                             % Ambiente de referencias utilizado.
\bibliographystyle{apalike}                     % Estilo de referencias APA.
\def\biblio{\clearpage\bibliographystyle{apalike}\bibliography{Referencias}} % Define el comando \biblio para referencias en subarchivos - NO CAMBIAR

% Configuración de listings como fallback para minted
% \usepackage{minted}                             % Permite darle formato al código que se muestre.
% Fallback: use listings (no -shell-escape required). If you prefer minted, uncomment the minted line above and compile with -shell-escape.
\usepackage{listings}
\usepackage{environ}

% Define un estilo para listings y un entorno tipo minted
\lstdefinestyle{myListing}{
  basicstyle=\ttfamily\small,
  breaklines=true,
  frame=single,
  numbers=left,
  numbersep=5pt,
  showstringspaces=false,
  keywordstyle=\color{blue}\bfseries,
  commentstyle=\color{gray}\itshape,
  stringstyle=\color{red}
}
\lstset{style=myListing}

\NewEnviron{minted}[1]{%
  \lstset{style=myListing, language=#1}%
  \begin{lstlisting}
\BODY
  \end{lstlisting}
}

% Enlaces clickeables y estilo de URL
\usepackage{hyperref}   % Permite que todo el documento sea clickeable.
\urlstyle{sf}                                   % Tipo de url
% ----- Cabecera y pies -----
\pagestyle{fancy}                           % Se define el estilo fancy
\fancyhead[RO,LE]{\thepage}                 % Número de página en la izquierda para par y derecha para impar
\fancyhead[RE,LO]{\nouppercase{\rightmark}} % Nombre del capítulo en la derecha para par y la izquierda para impar en la cabecera
%\renewcommand{\headrulewidth}{0pt}         % Cambiar para línea más gruesa
\fancyfoot{}                                % Saca el número de la página abajo.
\setlength{\headheight}{15pt}               % Altura suficiente para fancyhdr
\addtolength{\topmargin}{-2.5pt}            % Compensa la nueva altura de cabecera

\fancypagestyle{plain}{                     % Se redefine el estilo automático (plain) para que calce con el resto. En particular la 1ra página de cada capítulo
\fancyhf{}                                  % Elimina la cabecera y los pies
\fancyhead[RO,LE]{\thepage}                 % Número de página en la izquierda para par y derecha para impar
\fancyhead[RE,LO]{\nouppercase{\leftmark}}  % Nombre del capítulo en la derecha para par y la izquierda para impar en la cabecera
%\renewcommand{\headrulewidth}{0pt}         % Cambiar para línea más gruesa
\fancyfoot{}                                % Elimina el número de la página abajo
}


\fancypagestyle{resumen}{                   % Se redefine el estilo resumen para que calce con el resto. 
\fancyhf{}                                  % Elimina la cabecera y los pies
\fancyhead[RO,LE]{\thepage}                 % Número de página en la izquierda para par y derecha para impar
\fancyhead[RE,LO]{\nouppercase{Resumen}}    % Nombre del capítulo en la derecha para par y la izquierda para impar en la cabecera
%\renewcommand{\headrulewidth}{0pt}         % Cambiar para línea más gruesa
\fancyfoot{}                                % Elimina el número de la página abajo
}

\fancypagestyle{abstract}{                  % Se redefine el estilo resumen para que calce con el resto. 
\fancyhf{}                                  % Elimina la cabecera y los pies
\fancyhead[RO,LE]{\thepage}                 % Número de página en la izquierda para par y derecha para impar
\fancyhead[RE,LO]{\nouppercase{Abstract}}   % Nombre del capítulo en la derecha para par y la izquierda para impar en la cabecera
%\renewcommand{\headrulewidth}{0pt}         % Cambiar para línea más gruesa
\fancyfoot{}                                % Elimina el número de la página abajo
}

\fancypagestyle{agradecimientos}{                   % Se redefine el estilo resumen para que calce con el resto. 
\fancyhf{}                                          % Elimina la cabecera y los pies
\fancyhead[RO,LE]{\thepage}                         % Número de página en la izquierda para par y derecha para impar
\fancyhead[RE,LO]{\nouppercase{Agradecimientos}}    % Nombre del capítulo en la derecha para par y la izquierda para impar en la cabecera
%\renewcommand{\headrulewidth}{0pt}                 % Cambiar para línea más gruesa
\fancyfoot{}                                        % Elimina el número de la página abajo
}

% ----- Cabecera de la portada ----- 
\fancypagestyle{frontpage}{             % Se define el estilo frontpage.
	\fancyhf{}                          % Elimina la cabecera y los pies
	\renewcommand{\headrulewidth}{0pt}  % Elimina líneas en cabecera
	\renewcommand{\footrulewidth}{0pt}  % Elimina líneas en pies
	\vspace*{1\baselineskip}
	
	\fancyhead[R]{UNIVERSIDAD DE CONCEPCIÓN                        % Línea 1 de la cabecera derecha, debe ir el nombre de la UdeC en mayúsculas.
	\linebreak       FACULTAD DE CIENCIAS FÍSICAS Y MATEMÁTICAS}   % Línea 2 de la cabecera derecha, debe ir el nombre de la facultad en mayúsculas.
	\fancyhead[L]{ \includegraphics[width=0.7in]{escudo_udec.png}} % LOGO que irá en la cabecera izquierda.
}

% ----- Enlaces clickeables --------
\newcommand\myshade{85} % Permite la redefinición de colores a gusto del usuario

% Para elegir colores propios mirar los nombres relacionados con dvipsnames, aquí un url con los nombres de dvipsnames: https://www.overleaf.com/learn/latex/Using_colours_in_LaTeX

\colorlet{mylinkcolor}{DarkOrchid}   % Define Color de los enlaces clickeables internos (índice, referencias cruzadas, etc.). En este caso DarkOrchid
\colorlet{mycitecolor}{YellowOrange} % Define Color de las citas ubicadas en las referencias. En este caso YellowOrange
\colorlet{myurlcolor}{Aquamarine}    % Define Color de los enlaces url. En este caso Aquamarine

% Para dejar el documento sin texto en colores cambiar las tres líneas anteriores a Black.

\hypersetup{  %Define la forma en que se verán las cosas clickeables.
  	linkcolor  = mylinkcolor!\myshade!black,    % Aplica el color definido arriba. En este caso DarkOrchid
  	citecolor  = mycitecolor!\myshade!black,    % Aplica el color definido arriba. En este caso YellowOrange
  	urlcolor   = myurlcolor!\myshade!black,     % Aplica el color definido arriba. En este caso Aquamarine
  	colorlinks = true,                          % Elimina las cajas alrededor de lo clickeable y lo reemplaza por palabras a color.
}

\begin{document}
\pagenumbering{roman}                       % Usa numeración romana desde la portada
\def\biblio{}   % Resetea el comando biblio, de lo contrario una lista de referencias será producida después de cada capítulo
                % resets the biblio command, if not here a new reference list will be produced after every chapter


\begin{titlepage}
	
	\newgeometry{top=1 in, bottom=1 in, left=1 in, right= 1 in} 
	
	\thispagestyle{frontpage}
	
	\begin{center}
		
		\vspace*{4\baselineskip}
	
		
		{\Huge \textbf{BAYESIAN CHARACTERIZATION OF 241 NEW POSSIBLE OPEN CLUSTERS USING HDBSCAN AND GAIA DR3\\}}%No abreviar, no subrayar, no usar comillas. Se escribe completamente en mayúsculas
		        \vspace*{1.5\baselineskip}

		\large{\textbf{By: Lucas Matías Pulgar Escobar}}\\ %Nombre como aparece en registro académico
		
		\vspace{1,5\baselineskip}
		
		\large{Thesis presented to the Facultad de Ciencias Físicas y Matemáticas of the Universidad de Concepción for the academic degree of Master in Astronomy.} %Cambiar a grado académico correspondiente. 
		
		\vspace{1,5\baselineskip}
		November 2024\\ %Mes y año de la tesis, solo primera letra del mes en mayúsculas
		Concepción, Chile %Ciudad y país de publicación.
\vspace{1,5\baselineskip}

		\large{\textbf{Advisor: Ronald Enrique Mennickent Cid}}\\ %Nombre del profesor que dirigió el trabajo, de la comisión informante y otros asesores precedidos por los títulos correctos de la unidad académica. Si hay página de calificaciones, solo se menciona al profesor guía en la portada

	\end{center}
	
	\vspace*{4\baselineskip}
	
	% La cabecera de esta página se cambia en el documento principal, en "cabecera portada"
	
\end{titlepage}   % Genera la portada desde el archivo Portada.tex . Para editar este archivo ir a Otros -> Portada.tex
\vfill

%----------------Página de derechos de autor: elegir entre a) o b) y borrar/comentar la opción NO utilizada-----------------
\thispagestyle{empty}
\mbox{}                         % Ayuda a bajar el texto
\vfill                          % Deja el texto al fondo
\textcopyright\ 2024, Lucas Pulgar Escobar \\ % Derechos de autor

Reproduction in part or in total, for academic purposes, by any means or process, including the bibliographic citation of the document, is authorized.
\vspace{1cm}    % lo separa del fondo
\restoregeometry % Devuelve los márgenes después de la portada

%----------------Página de dedicatoria (opcional), descomentar para generar ---------------------------------

% Editar en Otros -> Dedicatoria.tex

%La dedicatoria debe ser ubicada en el extremo inferior derecho de una página vacía, está página es opcional
\thispagestyle{empty}
\mbox{}
\vfill
\hfill \text{To all the people who were part of this journey, to my mom, who supported me in many ways, to my friends, who helped, listened and congratulated me on several occasions, and especially to my significant other, Nicolas, with whom we spent so many nights on this road, sometimes tangentially, sometimes in parallel, and sometimes even asymptotically.}%Escribir aquí la dedicatoria.
   % Genera la pagina de calificaciones del archivo calificaciones.tex
\restoregeometry        % Devuelve los márgenes después de la página


%-----------------Página de agradecimientos (opcional), se incluye normalmente-------------------

% Editar en Otros -> Agradecimientos.

\newpage
\addcontentsline{toc}{chapter}{AGRADECIMIENTOS}  % Agrega esta sección al índice
\section*{AGRADECIMIENTOS}                       % Debe ir en mayúsculas por reglamento de la UdeC, tiene asterisco para no ser numerada.
    \subfile{Otros/Agradecimientos}              % Llama al archivo Agradecimientos en la carpeta Otros.


%-----------------Página de resumen (abstract)-------------

% Si la unidad académica lo requiere, se edita en  Otros -> Resumen.tex . El mismo resumen puede ser incluido en inglés (abstract) en la página siguiente, para agregarlo hay un espacio destinado en el mismo archivo antes mencionado.

\newpage
\addcontentsline{toc}{chapter}{Resumen} % Agrega esta sección al índice
\section*{Resumen}                      % Con asterisco para que no sea numerada.
    \subfile{Otros/Resumen}             % Llama al archivo Resumen en la carpeta Otros.

%--------------Página de índice.  

%\nocite{*}     % Des-comentar si se desea que TODAS las referencias sean impresas en la lista de referencias, incluyendo las que no fueron finalmente citadas en el texto.

\newpage
{\setstretch{1.0}   % Interlineado de la lista.
\tableofcontents
}

\newpage
{\setstretch{1.0} 
\listoftables}

\newpage
{\setstretch{1.0} 
\listoffigures}


\newpage
\addtocontents{toc}{\protect\setcounter{tocdepth}{4}}   % La profundidad del índice queda en 4, 1.1.1.1
\pagenumbering{arabic}                                  % Comienza la numeración arábiga (números normales)
\setcounter{page}{1}                                    % Comienza el contador de páginas en 1

% A continuación se dejan nombres de diversos capítulos o secciones, para cambiar el nombre del archivo tan solo se debe hacer en la carpeta "capitulos" y luego llamarlos de la forma correcta en "\subfile{Capitulos/nuevonombre}".
% Los nombres de los archivos no pueden llevar tíldes ni espacios para el correcto funcionamiento del compilador, esto no tiene nada que ver con que tengan o no tilde en el documento final.

\chapter{Introducción}                      % Nombre de section/chapter, este si usa tilde y es lo que se verá impreso
    \subfile{Capitulos/01Introduccion}      % Incluye el subarchivo del section/chapter 
\clearpage                                  % limpia la página luego de que el capítulo concluya.

\chapter{Marco Teórico}
    \subfile{Capitulos/02MarcoTeorico}
\clearpage

\chapter{Metodología}
    \subfile{Capitulos/03Metodologia}
\clearpage
  
\chapter{Análisis}
    \subfile{Capitulos/04Analisis}
\clearpage

\chapter{Discusión}
    \subfile{Capitulos/05Discusion}
\clearpage

\chapter{Conclusión}
    \subfile{Capitulos/06Conclusion}
\clearpage

\newpage
\renewcommand\refname{Referencias}          % Nombre para la lista de referencias, también se utiliza "Bibliografía"
{\setstretch{1.0}                           % Interlineado de las referencias 
\addcontentsline{toc}{chapter}{Referencias} % Cambia el nombre de la lista de referencias en el índice 
\bibliography{Referencias}                  % Agrega las referencias al documento, estas se ubican en el archivo Referencias.bib
}

\newpage
\renewcommand{\appendixpagename}{Apéndices}     % Nombre al inicio.
\addcontentsline{toc}{chapter}{Apéndices}       % Agrega "Apéndices" al índice

\appendix   % Empieza el ambiente de apéndices, desde ahora en adelante los capítulos, secciones, tablas, figuras, etc. vuelven a empezar su numeración

\chapter{Test}                      % Nombre del apéndice, el capítulo será "Apéndice A" en vez de "Capítulo 1"
    \subfile{Capitulos/Apendice}    % Llama al archivo apéndice
\clearpage

% Este segundo apéndice tiene información general sobre LaTeX. Está comentado para que no aparezca en tu documento, pero si lo deseas puedes descomentarlo para ver su contenido.

%\chapter{Algunos consejos sobre \LaTeX{}}
%    \subfile{Otros/Consejos}
%\clearpage

% ÉXITO EN TU TESIS

\end{document}
